\chapter{Reliable supervised-learning systems}
\label{ch:reliable-supervised-systems}

\begin{learningobjectives}
After this chapter, you should be able to:
\begin{itemize}
  \item separate training, prediction, and feedback paths;
  \item define a versioned model artifact and feature contract;
  \item monitor service behavior, input shift, predictions, and delayed outcomes; and
  \item design promotion, rollback, and retraining as governed operations.
\end{itemize}
\end{learningobjectives}

\begin{chapterconnection}
The preceding chapters developed candidate estimators and honest selection. This
Part II synthesis places an approved artifact inside the end-to-end architecture
of Chapter~\ref{ch:end-to-end}. A prediction system is incomplete until future
outcomes can test its claims and operators can respond safely.
\end{chapterconnection}

\section{Three paths define the operated system}

The \textbf{training path} validates a data snapshot, constructs split and
features, fits candidates, evaluates them, and promotes an immutable artifact.
The \textbf{prediction path} validates a request or batch, applies the exact
feature transformation and approved artifact, and records output context. The
\textbf{feedback path} joins predictions to outcomes after the target horizon
and produces evidence for continuation, investigation, rollback, or retraining.

\begin{center}
\begin{tikzpicture}[node distance=5mm]
  \node[flowbox] (data) {Versioned\\observations};
  \node[flowbox,right=of data] (train) {Train and\\evaluate};
  \node[flowbox,right=of train] (registry) {Approved\\artifact};
  \node[flowbox,below=of registry] (serve) {Batch or API\\prediction};
  \node[flowbox,left=of serve] (client) {Scientific\\client};
  \node[flowbox,left=of client] (outcome) {Delayed\\outcomes};
  \draw[flowarrow] (data) -- (train);
  \draw[flowarrow] (train) -- (registry);
  \draw[flowarrow] (registry) -- (serve);
  \draw[flowarrow] (serve) -- (client);
  \draw[dataarrow] (client) -- (outcome);
  \draw[dataarrow] (outcome) -- (data);
\end{tikzpicture}
\end{center}

Each arrow requires identity, schema, version, time, error, retry, access, and
retention semantics. The paths may run at different cadences and on different
infrastructure, but their contracts must agree.

\section{A model artifact is more than parameters}

An approved artifact binds together a feature schema and order, fitted
transformations, model state, class or output semantics, supported runtime,
training-data and code references, evaluation report, and model card. A registry
records lifecycle state such as candidate, approved, superseded, or withdrawn.

Never load arbitrary serialized objects from untrusted sources. Many Python
serialization formats can execute code. Control provenance, integrity, access,
and runtime compatibility. Prefer constrained formats when they represent the
needed operations faithfully.

\section{Offline and online features must agree}

Training-serving skew occurs when feature computation differs between historical
training and live prediction. Causes include time-zone handling, category
vocabularies, window boundaries, late data, default values, code versions, and
units. Reuse package code where possible, and test representative raw records
through both paths.

Point-in-time correctness is essential: a historical feature must use only
information available at its historical prediction time. Recomputing it from a
current database can accidentally include future corrections or outcomes.

\conceptcheck{A live feature has the same name, dtype, and range as its training
counterpart. Give two ways it can still have different semantics.}

\section{Prediction interfaces fail explicitly}

Batch inference favors throughput, snapshot lineage, restartable partitions,
and idempotent outputs. Online inference adds latency, concurrency, admission
control, timeout, and availability responsibilities. Both validate schemas,
bound resources, record artifact identity, and reject incompatible requests.

The interface should define behavior for missing, unknown, infinite, or out-of-
support inputs. Abstention or routing to human review can be safer than a forced
prediction. Error responses must not leak secrets or silently substitute a
different model.

\section{Monitoring has layers}

Operational monitoring covers request rate, latency, errors, saturation,
artifact load, and dependency health. Data-quality monitoring covers schema,
missingness, range, category, freshness, and join behavior. Distribution
monitoring compares feature and score summaries with reference windows.

Covariate shift means the distribution of $X$ changes. Prior shift changes
class prevalence. Concept shift changes the relationship between $X$ and $Y$.
Unlabeled inputs can reveal the first imperfectly; they cannot by themselves
measure predictive error or identify the third. Delayed outcomes are therefore
not optional if performance claims are to remain testable.

Monitor slices defined before release and investigate emerging slices carefully.
Alerts need thresholds, owners, runbooks, and actions. A dashboard without a
response policy is visualization, not operational control.

\begin{misconceptionbox}
\textbf{Misconception:} feature drift means the model must be retrained. 
\textbf{Correction:} drift is evidence to investigate. Retraining can preserve
the same defect, learn from corrupted labels, or degrade important groups.
Action depends on outcomes, cause, severity, and the validated response policy.
\end{misconceptionbox}

\section{Retraining is a release, not a reflex}

A retraining trigger begins a candidate-building workflow. The candidate must
pass data validation, tests, evaluation, subgroup and resource gates, and
approval. Automatic promotion is appropriate only when those gates and rollback
controls match the consequence of failure.

Rollback requires a previously known artifact, compatible feature interface,
and a procedure tested before the incident. Some shifts make rollback
insufficient because the old artifact faces the same new population. A safe
fallback may be a baseline, abstention, reduced scope, or human review.

\begin{professionalbox}
Separate the decision to recompute a candidate from the decision to release it.
Schedule-driven retraining is an opportunity to evaluate new evidence, not a
guarantee that the newest artifact is better.
\end{professionalbox}

\section{The release dossier}

A defensible release records the prediction contract, data and feature lineage,
split specification, candidate search, final evaluation, known limitations,
model and policy versions, security review, test evidence, monitoring plan,
owner, rollback plan, and expiration or review date. CI verifies reproducible
software properties; human review owns scientific and operational judgment that
cannot be reduced to one test.

\begin{workedexample}{an incident closes the evidence loop}
After release, the degradation service remains fast and error-free, but a new
batch produces higher scores and more out-of-support warnings. Outcomes will not
arrive for 100 cycles. The runbook narrows automated use and routes affected
cells to review rather than immediately retraining.

Investigation finds a sensor firmware change that altered one summary feature.
The team reconstructs the feature with point-in-time data, evaluates both the
approved artifact and a candidate on the corrected representation, and records
the incident. A compatibility adapter restores the defined unit; the model is
not changed. Operational evidence prevented a data-interface defect from being
misdiagnosed as concept drift.
\end{workedexample}

\section{Exercises}
\begin{enumerate}
  \item Draw training, prediction, and feedback paths for batch inference. Mark
  every durable artifact and join key.
  \item Specify a feature-parity test involving time windows, late records,
  unknown categories, and units.
  \item Design monitoring that distinguishes service failure, input shift, score
  shift, calibration change, and outcome degradation.
  \item Write a rollback rehearsal and explain one shift for which rollback
  would not restore acceptable behavior.
\end{enumerate}

\begin{chapterreview}
A reliable supervised-learning system joins three paths. Training creates and
promotes versioned evidence; prediction applies an approved artifact through a
validated interface; feedback connects outputs to delayed outcomes. Monitoring
must cover operations, data, scores, slices, and eventual performance. Drift
prompts investigation, while retraining and rollback remain governed releases.
\end{chapterreview}

\begin{notebooklink}
Lecture 17 notebook 00 maps the three system paths. Companion capstone
laboratories produce a release dossier, monitoring specification, outcome join,
and tested incident and rollback exercise.
\end{notebooklink}
